%% ──────────────────────────────────────────────
%%  پیوست – طبقه‌بندی مداخلات خارجی:
%%  از کشورگشایی تا همبستگی
%%  فایل: appendices/app-taxonomy-interventions.tex
%% ──────────────────────────────────────────────
\chapter{طبقه‌بندی مداخلات خارجی: از کشورگشایی تا همبستگی}
\label{ch:taxonomy-interventions}

\begin{tcolorbox}[mybox, title={درباره‌ی این پیوست}]
در فصل‌های ۱۵ تا ۱۷ از «مداخله‌ی خارجی» سخن گفتیم — اما این اصطلاح طیف
وسیعی از پدیده‌های تاریخی را در بر می‌گیرد.  این پیوست تلاش می‌کند
\textbf{نقشه‌ای مفهومی} از انواع مختلف دخالت خارجی در سرنوشت ملت‌ها ترسیم
کند: از استعمار کلاسیک تا مداخله‌ی بشردوستانه، از نظام مَندیت تا جنگ‌های
نیابتی.  هدف، هم شفاف‌سازی مفهومی است و هم نشان‌دادن اینکه «آمار شکست»
ارائه‌شده در فصل ۱۵ دقیقاً شامل کدام طبقه‌ها می‌شود.
\end{tcolorbox}


%% ================================================================
\section{چرا طبقه‌بندی مهم است؟}
\label{sec:tax-why}
%% ================================================================

\needspace{6\baselineskip}
یکی از رایج‌ترین خطاهای تحلیلی در بحث مداخله، \concept{خلط طبقات} است:

\begin{tcolorbox}[enemybox, title={سه خطای رایج}]
\textbf{خطای ۱ — بزرگ‌نمایی منتقدان:}
استعمار بریتانیا در هند، مَندیت فرانسه در سوریه، و حمله‌ی آمریکا به
عراق همه زیر عنوان «مداخله‌ی خارجی» قرار می‌گیرند تا نشان داده شود «هر
دخالتی» فاجعه‌بار است.

\textbf{خطای ۲ — خوش‌بینی مدافعان:}
بازسازی موفق ژاپن و آلمان پس از جنگ جهانی دوم (طبقه‌ی ۲: بازآرایی
پساجنگی) به‌عنوان دلیل «موفقیت مداخله» ارائه می‌شود — درحالی‌که این
موارد اساساً با حمله به عراق ۲۰۰۳ (طبقه‌ی ۵/۷) قابل مقایسه نیستند.

\textbf{خطای ۳ — تقلیل تاریخ:}
نظام مَندیت عثمانی «مداخله‌ی بشردوستانه» نامیده می‌شود — درحالی‌که در
واقع \keyword{تقسیم غنائم جنگی} بود.  به‌عکس، استعمارزدایی «شکست
مداخله» نامیده می‌شود — درحالی‌که اتفاقاً \keyword{پایان مداخله} بود.
\end{tcolorbox}


%% ================================================================
\section{هشت طبقه‌ی دخالت خارجی}
\label{sec:tax-eight-types}
%% ================================================================

\needspace{6\baselineskip}
در ادبیات روابط بین‌الملل \textbf{یک طبقه‌بندی مورد اجماع} وجود ندارد.
\thinker{جیمز روزنا} (\lat{James Rosenau}) در ۱۹۶۹ اولین تلاش نظام‌مند را
انجام داد؛\footnote{Rosenau, James N.  ``Intervention as a Scientific
Concept.'' \textit{Journal of Conflict Resolution} 13, no.\,2 (1969):
149--171.} \thinker{مارتا فینمور} (\lat{Martha Finnemore}) در ۲۰۰۳
تحول تاریخی «معنای مشروع مداخله» را ردیابی کرد؛\footnote{Finnemore,
Martha.  \textit{The Purpose of Intervention: Changing Beliefs about the
Use of Force}.  Cornell UP, 2003.}
و \thinker{الکساندر داونز} (\lat{Alexander Downes}) در ۲۰۲۱ مشخصاً
«تغییر رژیم تحمیلی» (\lat{FIRC}) را بررسی کرد.

بر اساس ترکیب این آثار و واقعیت‌های تاریخی، طبقه‌بندی هشت‌گانه‌ی زیر
پیشنهاد می‌شود.


%% ────────────────────────────────────
\subsection{طبقه‌ی ۱: کشورگشایی و استعمار کلاسیک}
\label{sec:tax-type1}

\begin{tcolorbox}[casebox, title={تعریف و ویژگی‌ها}]
\textbf{تعریف:} تصرف سرزمین دیگری برای بهره‌برداری اقتصادی، استراتژیک،
یا جمعیتی — معمولاً بلندمدت و بدون قصد بازگشت.

\textbf{انگیزه‌ی اصلی:} سلطه و منفعت مادی.

\textbf{توجیه رایج:}
\begin{itemize}
\item «رسالت تمدنی» (\lat{mission civilisatrice}) — فرانسه
\item «بار مرد سفید» (\lat{White Man's Burden}) — بریتانیا
\item «سرنوشت آشکار» (\lat{Manifest Destiny}) — آمریکا
\item برتری نژادی — همه‌ی قدرت‌های استعماری
\end{itemize}

\textbf{ویژگی متمایز:} استعمارگر ادعای «نجات» ندارد — یا اگر دارد،
نجات از «توحش» است، نه از «ظلم حاکم محلی».  هدف نه تغییر رژیم بلکه
\keyword{حذف حکومت محلی} و جایگزینی آن با اداره‌ی مستقیم خارجی است.
\end{tcolorbox}

\begin{tcolorbox}[wavebox, title={زیرشاخه‌های استعمار}]

\textbf{الف) استعمار استیطانی (\lat{Settler Colonialism}):}
\begin{itemize}
\item \textbf{ماهیت:} مهاجرت انبوه استعمارگران + جابجایی یا نابودی جمعیت
  بومی
\item \textbf{مثال‌ها:} آمریکای شمالی، استرالیا، نیوزیلند، الجزایر فرانسه،
  فلسطین/اسرائیل، رودزیا
\item \textbf{ویژگی خاص:} مهاجران قصد بازگشت ندارند؛ سرزمین را «مال خود»
  می‌دانند
\item \textbf{پیامد بلندمدت:} عمیق‌ترین و ماندگارترین شکل دخالت — آثارش
  قرن‌ها ادامه دارد
\end{itemize}

\tcblower

\textbf{ب) استعمار بهره‌برداری (\lat{Extractive Colonialism}):}
\begin{itemize}
\item \textbf{ماهیت:} بدون مهاجرت انبوه؛ هدف استخراج منابع و نیروی کار
\item \textbf{مثال‌ها:} کنگوی بلژیک، هند بریتانیا (بخش عمده)، اندونزی
  هلند
\item \textbf{ویژگی خاص:} جمعیت محلی «ابزار تولید» است، نه رقیب
\item \textbf{پیامد:} ساختارهای اقتصادی وابسته که پس از استقلال هم باقی
  ماندند
\end{itemize}

\tcblower

\textbf{ج) استعمار شرکتی (\lat{Corporate Colonialism}):}
\begin{itemize}
\item \textbf{ماهیت:} یک شرکت خصوصی — نه دولت — حکومت می‌کند
\item \textbf{مثال‌ها:} کمپانی هند شرقی بریتانیا \dated{۱۶۰۰–۱۸۵۸}،
  کمپانی هند شرقی هلند (\lat{VOC})، شرکت آفریقای جنوبی بریتانیا (سسیل رودز)
\item \textbf{ویژگی خاص:} انگیزه‌ی سود شرکتی مستقیم‌تر از منافع ملی
  است — \concept{خصوصی‌سازی سلطه}
\item \textbf{اهمیت امروزی:} الگوی شرکت‌های نظامی خصوصی مثل
  \lat{Blackwater}/\lat{Academi} و شرکت‌های استخراجی چندملیتی
\end{itemize}

\tcblower

\textbf{د) استعمار اجاره‌ای / امتیازی (\lat{Concession Colonialism}):}
\begin{itemize}
\item \textbf{ماهیت:} تصرف بخشی از سرزمین (بندر، جزیره، منطقه) از طریق
  «اجاره» یا «امتیاز» — اغلب با زور یا تهدید
\item \textbf{مثال‌ها:} هنگ‌کنگ \dated{۱۸۴۲–۱۹۹۷}، ماکائو، امتیازات
  خارجی در چین، کانال سوئز
\item \textbf{ویژگی خاص:} ظاهراً «قانونی» (معاهده‌ی دوجانبه)، اما
  معاهده‌ها با زور تحمیل شده بودند — چینی‌ها اینها را
  \keyword{معاهدات نابرابر} (\lat{不平等条约}) می‌نامند
\end{itemize}
\end{tcolorbox}

\begin{caseStudy}{هند بریتانیا: سه قرن، سه مرحله}
هند نمونه‌ی بارزی از تطور شکل‌های استعمار است:

\begin{description}[leftmargin=!, labelwidth=4.5cm,
  font=\bfseries\color{PrimaryMid}]
\item[مرحله ۱ \dated{۱۶۰۰–۱۷۵۷}]
  \concept{نفوذ تجاری} — کمپانی هند شرقی به‌عنوان بازرگان وارد شد
\item[مرحله ۲ \dated{۱۷۵۷–۱۸۵۸}]
  \concept{استعمار شرکتی} — پس از نبرد پلاسی، کمپانی عملاً
  حکومت‌کننده شد
\item[مرحله ۳ \dated{۱۸۵۸–۱۹۴۷}]
  \concept{استعمار دولتی مستقیم} — پس از شورش ۱۸۵۷ (که انگلیسی‌ها
  «شورش سپاهیان» و هندی‌ها «اولین جنگ استقلال» می‌نامند)، تاج
  بریتانیا مستقیماً حکومت را به‌دست گرفت
\item[مرحله ۴ \dated{۱۹۴۷}]
  \concept{استعمارزدایی فاجعه‌بار} — خروج شتاب‌زده + تقسیم =
  ۱–۲ میلیون کشته و ۱۵ میلیون آواره
\end{description}

\textbf{نکته‌ی کلیدی:} مشکل هند نه «مداخله برای نجات» بلکه
\keyword{پایان ۲۰۰ سال سلطه} بود.  تقسیم هند-پاکستان نتیجه‌ی «مداخله‌ی
خارجی» نبود — نتیجه‌ی \keyword{خروج بد مدیریت‌شده‌ی} استعمارگر بود.
\thinker{لرد ماونت‌بَتن} تاریخ خروج را ده ماه جلو انداخت، و خط مرزی
(\concept{خط رادکلیف}) توسط وکیلی انگلیسی که هرگز به هند نرفته بود در عرض
پنج هفته کشیده شد.\footnote{Khan, Yasmin.
\textit{The Great Partition: The Making of India and Pakistan}.
Yale UP, 2007.}
\end{caseStudy}


%% ────────────────────────────────────
\subsection{طبقه‌ی ۲: بازآرایی پساجنگی}
\label{sec:tax-type2}

\begin{tcolorbox}[casebox, title={تعریف و ویژگی‌ها}]
\textbf{تعریف:} تقسیم و سازمان‌دهی مجدد سرزمین‌ها پس از جنگ بزرگ،
توسط قدرت‌های فاتح.

\textbf{انگیزه‌ی اصلی:} مدیریت پیامدهای جنگ + تقسیم غنائم +
(ظاهراً) جلوگیری از جنگ بعدی.

\textbf{توجیه رسمی:} «نظم بین‌المللی جدید»، «صلح پایدار»،
«آمادگی ملت‌ها برای خودحکومتی».

\textbf{ویژگی متمایز:} این طبقه \textbf{بین استعمار و مداخله} قرار
دارد — هم عناصر امپریالیستی دارد (تقسیم غنائم) و هم ادعای «مدیریت
موقت» (قیمومیت).
\end{tcolorbox}

\begin{caseStudy}{نظام مَندیت: استعمار با نام جدید}
پس از فروپاشی امپراتوری عثمانی، جامعه‌ی ملل سرزمین‌های عثمانی را به
سه رده‌ی «مَندیت» (\lat{Mandate}) تقسیم
کرد:\footnote{Pedersen, Susan.  \textit{The Guardians: The League of
Nations and the Crisis of Empire}.  Oxford UP, 2015.}

\begin{table}[H]
\centering
\begin{tabularx}{\textwidth}{c|r|X|r}
\toprule
\textbf{رده} & \textbf{ادعای رسمی} & \textbf{واقعیت} & \textbf{نمونه} \\
\midrule
\textbf{A} & «نزدیک به استقلال» &
  استعمار با وعده‌ی مبهمِ آزادی؛ مرزها بر اساس منافع
  انگلیس و فرانسه &
  عراق، سوریه، لبنان، فلسطین \\
\textbf{B} & «نیازمند مدیریت بیشتر» &
  استعمار کلاسیک با لعاب حقوقی &
  تانگانیکا، کامرون، توگو \\
\textbf{C} & «بخشی از سرزمین قیّم» &
  الحاق عملی &
  نامیبیا، ساموآ \\
\bottomrule
\end{tabularx}
\end{table}

\textbf{ماده‌ی ۲۲ میثاق جامعه‌ی ملل:}
\begin{quote}\itshape
«مردمانی که هنوز قادر به اداره‌ی خود در شرایط سخت جهان مدرن نیستند
\ldots\ بهروزی و توسعه‌ی آنان امانتی مقدس تمدن است \ldots\ بهترین روش
تحقق این اصل، سپردن قیمومیت این مردمان به ملت‌های پیشرفته است.»
\end{quote}

\textbf{نقد:}
\begin{itemize}
\item \thinker{دیوید فرامکین} (\lat{David Fromkin}) در
  \textit{A Peace to End All Peace} \dated{۱۹۸۹} نشان داده که
  مرزکشی‌ها عمدتاً بر اساس \keyword{توافق مخفی سایکس-پیکو}
  \dated{۱۹۱۶} بودند — نه واقعیت‌های جمعیتی یا خواست مردم
\item \thinker{رشید خالیدی} (\lat{Rashid Khalidi}) در
  \textit{The Iron Cage} \dated{۲۰۰۶} استدلال کرده
  که مَندیت در فلسطین \textbf{هیچ تفاوتی} با استعمار مستقیم نداشت —
  فقط زبان حقوقی‌اش متفاوت بود
\item \thinker{یوجین روگان} (\lat{Eugene Rogan}) در
  \textit{The Fall of the Ottomans} \dated{۲۰۱۵} نشان داده
  که اعراب «آزادی» وعده داده شده بودند اما «قیمومیت» دریافت کردند — خیانتی
  که تا امروز حافظه‌ی جمعی منطقه را شکل می‌دهد
\end{itemize}
\end{caseStudy}

\begin{tcolorbox}[defbox, title={چرا مَندیت «مداخله‌ی خارجی» نبود؟}]
مَندیت‌های عثمانی با «مداخله‌ی خارجی» به‌معنای مدرن (مثلاً حمله به
عراق ۲۰۰۳) تفاوت‌های بنیادین دارند:

\begin{table}[H]
\centering
\begin{tabularx}{\textwidth}{r|X|X}
\toprule
\textbf{بُعد} & \textbf{مَندیت عثمانی} &
  \textbf{مداخله‌ی مدرن (مثلاً عراق ۲۰۰۳)} \\
\midrule
\textbf{زمینه} & پایان جنگ جهانی؛ فروپاشی امپراتوری &
  تصمیم یک‌جانبه‌ی قدرت خارجی \\
\textbf{ادعا} & «قیمومیت موقت» تا آمادگی برای استقلال &
  «رهایی» و «دموکراسی‌سازی» \\
\textbf{واقعیت} & تقسیم غنائم بین فاتحان &
  تغییر رژیم + بازسازی \\
\textbf{مخاطب} & سرزمین‌های بدون دولتِ شناخته‌شده (عثمانی فروپاشیده) &
  دولت‌های موجود و شناخته‌شده \\
\textbf{حقوقی} & مصوبه‌ی جامعه‌ی ملل (اجماع فاتحان) &
  اغلب بدون مجوز شورای امنیت \\
\textbf{مدت} & رسماً «تا آمادگی» (عملاً نامحدود) &
  رسماً «موقت» (عملاً طولانی) \\
\bottomrule
\end{tabularx}
\end{table}

\textbf{نتیجه:} مَندیت‌ها بهتر است در طبقه‌ی \concept{«بازآرایی
امپریالیستی پساجنگی»} قرار گیرند — نه «مداخله‌ی خارجی». اما
\keyword{آثار} آنها (مرزکشی مصنوعی، تنش‌های قومی-مذهبی، دولت‌های ضعیف)
دقیقاً همان زمینه‌ای را ایجاد کردند که بعدها «مداخله‌ی خارجی» را ممکن و
حتی ظاهراً «ضروری» ساخت.
\end{tcolorbox}

\textbf{سایر نمونه‌های بازآرایی پساجنگی:}

\begin{table}[H]
\centering
\begin{tabularx}{\textwidth}{r|X|r}
\toprule
\textbf{مورد} & \textbf{ماهیت} & \textbf{نتیجه‌ی بلندمدت} \\
\midrule
\histref{تقسیم آلمان}{۱۹۴۵} & بازآرایی پساجنگی + رقابت ایدئولوژیک &
  وحدت مجدد ۱۹۹۰ (موفقیت نسبی) \\
\histref{اشغال ژاپن}{۱۹۴۵–۵۲} & بازآرایی پساجنگی + بازسازی &
  دموکراسی پایدار (نمونه‌ی «موفق») \\
\histref{معاهده‌ی ورسای}{۱۹۱۹} & تجزیه‌ی آلمان و اتریش-مجارستان &
  ناکامی ← جنگ جهانی دوم \\
\histref{تقسیم کره}{۱۹۴۵} & تقسیم بین آمریکا و شوروی &
  جنگ ۱۹۵۰ + انقسام تا امروز \\
\bottomrule
\end{tabularx}
\end{table}


%% ────────────────────────────────────
\subsection{طبقه‌ی ۳: استعمارزدایی و مدیریت خروج}
\label{sec:tax-type3}

\begin{tcolorbox}[casebox, title={تعریف و ویژگی‌ها}]
\textbf{تعریف:} فرایند خروج قدرت استعماری و انتقال حاکمیت به مردم محلی.

\textbf{انگیزه‌ی اصلی:} فشار جنبش‌های استقلال‌طلبانه + ضعف اقتصادی و
نظامی قدرت استعماری + تغییر هنجارهای بین‌المللی.

\textbf{ویژگی متمایز:} مسئله نه \textbf{ورود} بیگانه بلکه \keyword{نحوه‌ی
خروجش} است — و آثاری که بر جای می‌گذارد.
\end{tcolorbox}

\begin{tcolorbox}[wavebox, title={سه شکل استعمارزدایی}]

\textbf{الف) خروج مذاکره‌ای (\lat{Negotiated Decolonization}):}
\begin{itemize}
\item قدرت استعماری و نخبگان محلی بر سر شرایط انتقال توافق می‌کنند
\item \textbf{مثال‌ها:} بسیاری از مستعمرات بریتانیا در آفریقا (غنا ۱۹۵۷،
  نیجریه ۱۹۶۰)؛ هنگ‌کنگ ۱۹۹۷ (مذاکره‌ی بریتانیا-چین)
\item \textbf{مزیت:} انتقال نسبتاً آرام
\item \textbf{مشکل:} نخبگان «دست‌نشانده» اغلب به قدرت می‌رسند
\end{itemize}

\tcblower

\textbf{ب) خروج خشونت‌آمیز (\lat{Violent Decolonization}):}
\begin{itemize}
\item قدرت استعماری با جنگ بیرون رانده می‌شود
\item \textbf{مثال‌ها:} الجزایر ۱۹۵۴–۶۲ ($\sim$۱.۵ میلیون کشته)؛
  ویتنام ۱۹۴۵–۵۴؛ موزامبیک و آنگولا
\item \textbf{مزیت:} استقلال «اصیل‌تر» — جنبش ملی قوی‌تر
\item \textbf{مشکل:} ویرانی + نظامی‌شدن جامعه + انتقام
\end{itemize}

\tcblower

\textbf{ج) خروج فاجعه‌بار (\lat{Catastrophic Decolonization}):}
\begin{itemize}
\item خروج شتاب‌زده + تقسیم سرزمین = فاجعه‌ی انسانی
\item \textbf{مثال‌ها:} هند/پاکستان ۱۹۴۷ (۱–۲ میلیون کشته)؛ فلسطین
  ۱۹۴۸ (۷۰۰٬۰۰۰ آواره‌ی فلسطینی = نکبت)
\item \textbf{مشکل بنیادین:} قدرت استعماری بدون حل مسائل بنیادین خارج
  شد — یا عامدانه مسائل را «باز» گذاشت تا نفوذش ادامه یابد
  (\concept{استعمار نو})
\end{itemize}
\end{tcolorbox}

\begin{caseStudy}{فلسطین: تلاقی چهار طبقه}
فلسطین پیچیده‌ترین مورد است زیرا \textbf{چهار طبقه هم‌زمان} در آن فعال‌اند:

\begin{enumerate}
\item \textbf{بازآرایی پساجنگی (طبقه‌ی ۲):} مَندیت بریتانیا
  \dated{۱۹۲۰–۱۹۴۸} — تقسیم سرزمین عثمانی
\item \textbf{استعمار استیطانی (طبقه‌ی ۱):} پروژه‌ی صهیونیستی با حمایت
  اعلامیه‌ی بالفور \dated{۱۹۱۷} — مهاجرت و استیطان
\item \textbf{استعمارزدایی فاجعه‌بار (طبقه‌ی ۳):} خروج بریتانیا
  \dated{۱۹۴۸} بدون حل مسئله — نکبت
\item \textbf{بازیابی/وحدت‌طلبی (طبقه‌ی ۶):} پان‌عربیسم و تلاش
  کشورهای عربی برای «آزادسازی فلسطین» — جنگ‌های ۱۹۴۸ و ۱۹۶۷
\end{enumerate}

خلط این طبقه‌ها یکی از دلایل اصلی پیچیدگی مسئله‌ی فلسطین است: هر طرف
یکی از این طبقه‌ها را برجسته و بقیه را پنهان می‌کند.
\end{caseStudy}


%% ────────────────────────────────────
\subsection{طبقه‌ی ۴: مداخلات ایدئولوژیک (جنگ سرد و فراتر)}
\label{sec:tax-type4}

\begin{tcolorbox}[casebox, title={تعریف و ویژگی‌ها}]
\textbf{تعریف:} ورود قدرت خارجی برای تغییر نظام سیاسی کشوری دیگر بر
اساس رقابت ایدئولوژیک.

\textbf{انگیزه‌ی اصلی:} رقابت آمریکا-شوروی (در دوره‌ی جنگ سرد)؛ رقابت‌های
منطقه‌ای (ایران-عربستان، هند-پاکستان).

\textbf{ویژگی متمایز:} معمولاً \keyword{مخفی یا نیابتی} — کودتا، مسلح‌کردن
شورشیان، ترور، عملیات روانی.

\textbf{توجیه رایج:} «جلوگیری از سقوط دومینوها»؛ «حمایت از خلق»؛
«صدور انقلاب».
\end{tcolorbox}

\begin{table}[H]
\centering
\begin{tabularx}{\textwidth}{r|r|r|X}
\toprule
\textbf{مورد} & \textbf{مداخله‌گر} & \textbf{ابزار} &
  \textbf{نتیجه‌ی بلندمدت} \\
\midrule
\histref{ایران}{۱۹۵۳} & آمریکا/بریتانیا & کودتا &
  ۲۵ سال دیکتاتوری $\to$ انقلاب ۱۹۷۹ \\
\histref{گواتمالا}{۱۹۵۴} & آمریکا & کودتا &
  ۳۶ سال جنگ داخلی و نسل‌کشی \\
\histref{شیلی}{۱۹۷۳} & آمریکا & کودتا &
  ۱۷ سال دیکتاتوری پینوشه \\
\histref{افغانستان}{۱۹۷۹} & شوروی & اشغال نظامی &
  ۱۰ سال جنگ $\to$ طالبان $\to$ القاعده \\
\histref{نیکاراگوئه}{۱۹۸۰ها} & آمریکا & مسلح‌کردن کنتراها &
  جنگ داخلی و ویرانی \\
\bottomrule
\end{tabularx}
\end{table}

\begin{tcolorbox}[enemybox, title={الگوی تکراری}]
تقریباً در \textbf{همه‌ی} مداخلات ایدئولوژیک جنگ سرد، یک الگوی مشابه
تکرار شد:
\begin{enumerate}
\item قدرت خارجی یک رهبر/حزب «دوست» را شناسایی می‌کند
\item با کودتا یا جنگ نیابتی، رژیم «دشمن» سرنگون می‌شود
\item رهبر «دوست» به قدرت می‌رسد — معمولاً دیکتاتور و فاسد
\item مردم رنج می‌برند؛ مقاومت شکل می‌گیرد
\item مقاومت سرکوب می‌شود — با حمایت همان قدرت خارجی
\item سرانجام، رژیم دست‌نشانده سقوط می‌کند — اما ویرانی باقی می‌ماند
\end{enumerate}

\textbf{نرخ موفقیت:} نزدیک به صفر. \thinker{لیندسی اُرکند}
(\lat{Lindsey O'Rourke}) در \textit{Covert Regime Change}
\dated{۲۰۱۸}\footnote{O'Rourke, Lindsey A. \textit{Covert Regime Change:
America's Secret Cold War}. Cornell UP, 2018.}
شصت‌وچهار عملیات مخفی آمریکا برای تغییر رژیم را بررسی کرده و نشان داده
که اکثریت قاطع آنها نه‌تنها به نتیجه‌ی مطلوب نرسیدند، بلکه وضعیت را
\keyword{بدتر} کردند.
\end{tcolorbox}


%% ────────────────────────────────────
\subsection{طبقه‌ی ۵: مداخلات «بشردوستانه»}
\label{sec:tax-type5}

\begin{tcolorbox}[casebox, title={تعریف و ویژگی‌ها}]
\textbf{تعریف:} مداخله‌ی نظامی با توجیه جلوگیری از نسل‌کشی، جنایات جنگی،
یا بحران انسانی.

\textbf{انگیزه‌ی اعلام‌شده:} حقوق بشر، دکترین \concept{مسئولیت حمایت}
(\lat{R2P}).

\textbf{انگیزه‌ی واقعی:} ترکیبی از بشردوستانه + ژئوپلیتیک — نسبت این
دو در هر مورد متفاوت است.

\textbf{ویژگی متمایز:} بالاترین ادعای اخلاقی، اما متناقض‌ترین نتایج.
\end{tcolorbox}

\begin{table}[H]
\centering
\begin{tabularx}{\textwidth}{r|r|X}
\toprule
\textbf{مورد} & \textbf{ادعا} & \textbf{نتیجه} \\
\midrule
\histref{سومالی}{۱۹۹۲} & نجات از قحطی &
  شکست نظامی + خروج ← «سندرم سومالی» \\
\histref{رواندا}{۱۹۹۴} & — (مداخله نشد) &
  ۸۰۰٬۰۰۰ کشته ← بحران وجدانی جهانی \\
\histref{کوزوو}{۱۹۹۹} & جلوگیری از پاکسازی قومی &
  نسبتاً موفق؛ اما بدون مجوز شورای امنیت \\
\histref{لیبی}{۲۰۱۱} & جلوگیری از کشتار &
  هرج‌ومرج، بازار برده، بحران مهاجرت \\
\histref{سوریه}{۲۰۱۱–} & — (مداخله‌ی مؤثر نشد) &
  ۵۰۰٬۰۰۰ کشته، ۱۲ میلیون آواره \\
\bottomrule
\end{tabularx}
\end{table}

\begin{tcolorbox}[legalbox, title={تناقض بنیادین R2P}]
دکترین «مسئولیت حمایت» \dated{۲۰۰۵} سه ستون دارد:
\begin{enumerate}
\item \concept{مسئولیت پیشگیری} (Responsibility to Prevent)
\item \concept{مسئولیت واکنش} (Responsibility to React)
\item \concept{مسئولیت بازسازی} (Responsibility to Rebuild)
\end{enumerate}

\textbf{مشکل:} در عمل، فقط ستون دوم (واکنش نظامی) اجرا می‌شود —
پیشگیری نادیده گرفته می‌شود (ارزان‌تر اما بی‌جذابیت) و بازسازی رها
می‌شود (گران‌تر و بی‌پایان‌تر).  لیبی نمونه‌ی بارز است: بمباران شد اما
\keyword{بازسازی نشد} — مثل جرّاحی که فقط ببُرد و بخیه نزند.
\end{tcolorbox}


%% ────────────────────────────────────
\subsection{طبقه‌ی ۶: مداخلات بازیابی و وحدت‌طلبانه}
\label{sec:tax-type6}

\begin{tcolorbox}[casebox, title={تعریف و ویژگی‌ها}]
\textbf{تعریف:} تلاش برای الحاق یا «بازپس‌گیری» سرزمین‌هایی که «متعلق
به ما» تلقی می‌شوند — بر اساس هویت قومی، زبانی، دینی، یا تاریخی.

\textbf{انگیزه‌ی اصلی:} ناسیونالیسم + حسرت امپراتوری ازدست‌رفته +
رقابت ژئوپلیتیک.

\textbf{توجیه رایج:} «حمایت از هم‌تباران»، «بازگشت به وحدت تاریخی»،
«اتحاد امت/ملت».

\textbf{ویژگی متمایز:} استفاده از \keyword{هویت} به‌عنوان ابزار مداخله —
که آن را از مداخلات صرفاً ژئوپلیتیک متمایز می‌کند (هرچند ژئوپلیتیک
همیشه در پس‌زمینه هست).
\end{tcolorbox}

\begin{table}[H]
\centering
\begin{tabularx}{\textwidth}{r|r|r|X}
\toprule
\textbf{مورد} & \textbf{ایدئولوژی} & \textbf{ادعا} & \textbf{نتیجه} \\
\midrule
\histref{وحدت مصر-سوریه}{۱۹۵۸} & پان‌عربیسم &
  «ملت عرب واحد» & شکست پس از ۳ سال \\
\histref{حمله‌ی عراق به کویت}{۱۹۹۰} & پان‌عربیسم/ناسیونالیسم عراقی &
  «استان ۱۹ عراق» & شکست نظامی + تحریم \\
\histref{روسیه در کریمه}{۲۰۱۴} & پان‌روسیسم &
  «حمایت از روس‌زبانان» & الحاق + جنگ دونباس \\
\histref{ترکیه در شمال سوریه}{۲۰۱۶–} & نئوعثمانی‌گری &
  «امنیت مرزی» + «حمایت از ترکمن‌ها» & اشغال نامحدود \\
\histref{پاکستان در کشمیر}{۱۹۴۷–} & ناسیونالیسم اسلامی &
  «جهاد آزادی‌بخش» & درگیری بی‌پایان \\
\bottomrule
\end{tabularx}
\end{table}

\begin{tcolorbox}[defbox, title={ایران بزرگ، عثمانی نو، و رؤیاهای دیگر}]
این طبقه شامل \concept{رؤیاهای احیای امپراتوری} نیز می‌شود — پدیده‌ای
که در قرن بیست‌ویکم دوباره قوت گرفته:

\begin{itemize}
\item \textbf{«عثمانی نو»:} سیاست خارجی ترکیه‌ی اردوغان، با تأکید بر
  نقش رهبری در «جهان اسلام» و حضور نظامی در سوریه، لیبی، عراق، و قطر
\item \textbf{«ایران بزرگ» / «هلال شیعه»:} نفوذ ایران از طریق نیروهای
  نیابتی در عراق، سوریه، لبنان، و یمن
\item \textbf{«دنیای روسی» (\lat{Русский мир}):} ادعای پوتین مبنی بر
  حمایت از روس‌زبانان در همه‌جا — توجیه حمله به اوکراین
\item \textbf{«خلافت»:} داعش — تلاش برای احیای خلافت اسلامی با زور
\item \textbf{پان‌ترکیسم:} رؤیای «اتحاد جهان تُرک» — از بالکان تا
  سین‌کیانگ
\end{itemize}

\textbf{ویژگی مشترک:} همه‌ی اینها از \keyword{حسرت عظمت گذشته} تغذیه
می‌کنند و نوعی \concept{مداخله} در سرنوشت دیگران را — به نام «بازگشت به
کل واحد» — مشروع می‌دانند.
\end{tcolorbox}


%% ────────────────────────────────────
\subsection{طبقه‌ی ۷: مداخلات ضدتروریستی/امنیتی}
\label{sec:tax-type7}

\begin{tcolorbox}[casebox, title={تعریف و ویژگی‌ها}]
\textbf{تعریف:} مداخله‌ی نظامی یا شبه‌نظامی با توجیه مقابله با تهدید
تروریستی.

\textbf{انگیزه‌ی اعلام‌شده:} امنیت ملی و بین‌المللی.

\textbf{ویژگی متمایز:} تعریف \concept{«تروریسم»} کاملاً سیاسی است — و
قابل کش‌وقوس برای پوشش‌دادن هر نوع مداخله‌ای. همان گروهی که برای یکی
«تروریست» است، برای دیگری «مبارز آزادی» است.

\textbf{فرمول کلیدی:} پس از ۱۱ سپتامبر ۲۰۰۱، اصطلاح «جنگ با ترور»
(\lat{War on Terror}) به \keyword{چک سفید} برای مداخله‌ی نظامی
بدل شد.
\end{tcolorbox}

\begin{table}[H]
\centering
\begin{tabularx}{\textwidth}{r|r|X}
\toprule
\textbf{مورد} & \textbf{توجیه} & \textbf{آثار جانبی} \\
\midrule
\histref{افغانستان}{۲۰۰۱–۲۰۲۱} & نابودی القاعده/طالبان &
  ۲۰ سال جنگ؛ بازگشت طالبان؛ ۲ تریلیون دلار هزینه \\
\histref{عراق}{۲۰۰۳} & «سلاح‌های کشتار جمعی» (دروغ) &
  داعش؛ جنگ فرقه‌ای؛ بی‌ثباتی منطقه‌ای \\
\histref{جنگ پهپادی آمریکا}{۲۰۰۴–} & «ترور هدفمند» تروریست‌ها &
  کشتار غیرنظامیان؛ نفرت‌پراکنی؛ بازتولید ترور \\
\histref{ترکیه در شمال سوریه}{۲۰۱۶–} & «مبارزه با تروریسم» (PKK/YPG) &
  پاکسازی قومی کردها؛ اشغال \\
\histref{چین در سین‌کیانگ}{۲۰۱۷–} & «مبارزه با افراطی‌گری» &
  بازداشت انبوه اویغورها؛ نسل‌کشی فرهنگی \\
\bottomrule
\end{tabularx}
\end{table}

%% ────────────────────────────────────
\subsection{طبقه‌ی ۸: مداخلات اقتصادی ساختاری}
\label{sec:tax-type8}

\begin{tcolorbox}[casebox, title={تعریف و ویژگی‌ها}]
\textbf{تعریف:} تحمیل مدل اقتصادی خاص بر کشوری دیگر از طریق ابزارهای
غیرنظامی — اما با آثاری که گاه ویرانگرتر از جنگ است.

\textbf{انگیزه‌ی اصلی:} منافع اقتصادی قدرت‌های بزرگ + ایدئولوژی
نئولیبرال + بازپرداخت بدهی‌ها.

\textbf{ابزارها:}
\begin{itemize}
\item \concept{برنامه‌های تعدیل ساختاری} (\lat{Structural Adjustment
  Programs — SAPs}) صندوق بین‌المللی پول
\item \concept{شرط‌گذاری} (\lat{Conditionality}): وام در ازای «اصلاحات»
  اقتصادی
\item \concept{تحریم‌های اقتصادی}: فشار بر کل جامعه (نه فقط رهبران)
\item \concept{معاهدات تجاری نابرابر}: دسترسی یک‌طرفه به بازارها
\item \concept{دیپلماسی بدهی} (\lat{Debt-Trap Diplomacy}): وام‌های
  بلندمدت با شرایط سنگین — الگوی چین در آفریقا و آسیا
\end{itemize}

\textbf{ویژگی متمایز:} بدون خون‌ریزی ظاهری، اما با آثار ساختاری عمیق و
ماندگار. \thinker{نائومی کلاین} (\lat{Naomi Klein}) در
\textit{The Shock Doctrine} \dated{۲۰۰۷}\footnote{Klein, Naomi.
\textit{The Shock Doctrine: The Rise of Disaster Capitalism}. Picador,
2007.} از اصطلاح \concept{«سرمایه‌داری فاجعه»}
(\lat{Disaster Capitalism}) استفاده کرده: بحران‌ها (طبیعی، نظامی، یا
مالی) فرصتی برای تحمیل «اصلاحات» نئولیبرال ایجاد می‌کنند.
\end{tcolorbox}

\begin{table}[H]
\centering
\begin{tabularx}{\textwidth}{r|r|X}
\toprule
\textbf{مورد} & \textbf{ابزار} & \textbf{پیامد} \\
\midrule
\histref{آمریکای لاتین}{دهه‌های ۱۹۸۰–۹۰} &
  برنامه‌های تعدیل ساختاری \lat{IMF} &
  خصوصی‌سازی، بیکاری انبوه، شورش‌های نان \\
\histref{روسیه}{۱۹۹۱–۹۸} &
  «شوک‌درمانی» (\lat{Shock Therapy}) &
  فروپاشی اقتصادی، ظهور الیگارش‌ها، فقر گسترده \\
\histref{یونان}{۲۰۱۰–۱۸} &
  برنامه‌ی ریاضت اقتصادی اتحادیه‌ی اروپا &
  ۲۵٪ کاهش تولید ناخالص، بیکاری ۲۷٪، مهاجرت نخبگان \\
\histref{عراق (پس از ۲۰۰۳)}{۲۰۰۳–} &
  خصوصی‌سازی اجباری توسط \lat{CPA} &
  نابودی صنایع داخلی، وابستگی به واردات \\
\histref{آفریقای صحرایی}{دهه‌های ۱۹۸۰–۲۰۰۰} &
  شرط‌گذاری بانک جهانی &
  کاهش بودجه‌ی آموزش و بهداشت، افزایش نابرابری \\
\histref{سریلانکا}{۲۰۲۲} &
  بحران بدهی (چین + بازارهای بین‌المللی) &
  فروپاشی اقتصادی، شورش مردمی، سقوط دولت \\
\bottomrule
\end{tabularx}
\end{table}

\begin{tcolorbox}[enemybox, title={چرا مداخله‌ی اقتصادی در آمار «شکست مداخلات» نمی‌آید؟}]
مداخلات اقتصادی ساختاری معمولاً در آمار «شکست مداخلات خارجی» لحاظ
\textbf{نمی‌شوند} — و دلایلش خود افشاگر است:

\begin{enumerate}
\item \textbf{تعریف محدود «مداخله»:} بیشتر پژوهش‌ها «مداخله» را به
  عمل \keyword{نظامی} محدود می‌کنند — فشار اقتصادی «مداخله» تلقی نمی‌شود
\item \textbf{پنهان‌بودن خشونت:} مرگ ناشی از بیمارستان بسته‌شده «خشونت»
  نامیده نمی‌شود — اما برای قربانی، فرقی با بمب ندارد
\item \textbf{مسئولیت‌گریزی:} \lat{IMF} و بانک جهانی می‌گویند «ما فقط
  توصیه کردیم؛ دولت‌ها خودشان تصمیم گرفتند» — درحالی‌که دولت‌ها
  «انتخابی» جز پذیرش نداشتند
\item \textbf{ایدئولوژی:} در پارادایم نئولیبرال، «بازار آزاد» نه مداخله
  بلکه \concept{«وضع طبیعی»} تلقی می‌شود — هر چیز دیگری «مداخله» است
\end{enumerate}

\textbf{اگر مداخلات اقتصادی هم لحاظ شوند}، آمار شکست مداخلات خارجی
احتمالاً از ۸۴٪ هم \keyword{بالاتر} خواهد رفت.
\end{tcolorbox}


%% ================================================================
\section{نقشه‌ی مفهومی: روابط بین طبقه‌ها}
\label{sec:tax-conceptual-map}
%% ================================================================

\needspace{8\baselineskip}
این هشت طبقه مستقل از هم نیستند — بلکه شبکه‌ای از روابط دارند:

\begin{tcolorbox}[timelinebox, title={روابط علّی و تاریخی بین طبقه‌ها}]

\textbf{زنجیره‌ی ۱: از استعمار تا مداخله}
\begin{quote}
استعمار (طبقه‌ی ۱) $\longrightarrow$
بازآرایی پساجنگی (طبقه‌ی ۲) $\longrightarrow$
استعمارزدایی فاجعه‌بار (طبقه‌ی ۳) $\longrightarrow$
دولت‌های ضعیف و مرزهای مصنوعی $\longrightarrow$
بحران و بی‌ثباتی $\longrightarrow$
مداخله‌ی بشردوستانه (طبقه‌ی ۵) یا ضدتروریستی (طبقه‌ی ۷)
\end{quote}

\textbf{مثال:} خاورمیانه —
استعمار عثمانی/اروپایی $\to$ سایکس-پیکو و مَندیت‌ها $\to$
استقلال‌های ناقص $\to$ دولت‌های استبدادی $\to$ بحران $\to$
حمله به عراق ۲۰۰۳ $\to$ داعش $\to$ مداخلات بیشتر

\tcblower

\textbf{زنجیره‌ی ۲: از ایدئولوژی تا تروریسم}
\begin{quote}
مداخله‌ی ایدئولوژیک جنگ سرد (طبقه‌ی ۴) $\longrightarrow$
دیکتاتوری‌های دست‌نشانده $\longrightarrow$
سرکوب و رادیکالیزاسیون $\longrightarrow$
تروریسم $\longrightarrow$
مداخله‌ی ضدتروریستی (طبقه‌ی ۷)
\end{quote}

\textbf{مثال:} افغانستان —
حمایت آمریکا از مجاهدین (طبقه‌ی ۴) $\to$ سقوط شوروی $\to$
جنگ داخلی $\to$ طالبان $\to$ القاعده $\to$ ۱۱ سپتامبر $\to$
حمله‌ی آمریکا ۲۰۰۱ (طبقه‌ی ۷) $\to$ ۲۰ سال جنگ $\to$ بازگشت طالبان

\tcblower

\textbf{زنجیره‌ی ۳: از اقتصاد تا بحران}
\begin{quote}
مداخله‌ی اقتصادی ساختاری (طبقه‌ی ۸) $\longrightarrow$
فقر و نابرابری $\longrightarrow$
شورش و بی‌ثباتی $\longrightarrow$
سرکوب $\longrightarrow$
مداخله‌ی بشردوستانه (طبقه‌ی ۵) یا وحدت‌طلبانه (طبقه‌ی ۶)
\end{quote}

\textbf{مثال:} بهار عربی —
دهه‌ها سیاست‌های نئولیبرال + فساد $\to$ نابرابری $\to$
شورش‌ها $\to$ سرکوب $\to$ جنگ داخلی (سوریه، یمن، لیبی) $\to$
مداخلات منطقه‌ای و بین‌المللی

\tcblower

\textbf{نتیجه‌ی کلیدی:} \concept{مداخله‌ها در خلأ رخ نمی‌دهند} — بلکه
حلقه‌هایی در زنجیره‌ی طولانیِ سلطه، بحران، و واکنش هستند. هر مداخله
زمینه‌ی مداخله‌ی بعدی را فراهم می‌کند — \keyword{چرخه‌ی باطل مداخله}.
\end{tcolorbox}


%% ================================================================
\section{آنچه آمار «۸۴٪ شکست» شامل می‌شود — و آنچه نمی‌شود}
\label{sec:tax-what-counts}
%% ================================================================

\begin{tcolorbox}[legalbox, title={محدوده‌ی دقیق آمار فصل ۱۵}]
آمار شکست مداخلات که در فصل~\ref{ch:paradox-of-liberation} ارائه شد،
عمدتاً بر پژوهش‌هایی استوار است که \concept{تغییر رژیم تحمیلی از خارج}
(\lat{Foreign-Imposed Regime Change — FIRC}) را بررسی کرده‌اند.  مهم‌ترین
این پژوهش‌ها:

\begin{description}[leftmargin=!, labelwidth=4cm,
  font=\bfseries\color{PrimaryMid}]
\item[\thinker{داونز} (2021)] ۱۲۰ مورد تغییر رژیم تحمیلی از ۱۸۱۶ تا
  ۲۰۰۸ — نتیجه: اکثریت قاطع به بی‌ثباتی و جنگ داخلی منجر شد
\item[\lat{RAND} (2019)] ۲۹ مورد ملت‌سازی آمریکایی — نتیجه: فقط ۴ مورد
  «موفقیت پایدار» (آلمان، ژاپن، پاناما، کرواسی) — و هر چهار مورد
  شرایط بسیار خاصی داشتند
\item[\thinker{اورکند} (2018)] ۶۴ عملیات مخفی آمریکا — نتیجه: اکثریت
  شکست خوردند یا وضع را بدتر کردند
\item[\lat{UCDP}] داده‌های درگیری‌های مسلحانه — نشان‌دهنده‌ی افزایش
  درگیری پس از مداخلات خارجی
\end{description}
\end{tcolorbox}

\begin{table}[H]
\centering\small
\begin{tabularx}{\textwidth}{r|c|c|X}
\toprule
\textbf{طبقه} & \textbf{در آمار هست؟} & \textbf{در آمار نیست} &
  \textbf{دلیل} \\
\midrule
۱. استعمار کلاسیک & & \checkmark &
  پدیده‌ی متفاوت؛ منطق سلطه‌ی بلندمدت \\
۲. بازآرایی پساجنگی & بعضاً\textsuperscript{*} & بعضاً &
  ژاپن و آلمان گاه «موفقیت» شمرده می‌شوند \\
۳. استعمارزدایی & & \checkmark &
  خروج است، نه ورود \\
۴. مداخله‌ی ایدئولوژیک & \checkmark & &
  هسته‌ی اصلی پژوهش‌ها \\
۵. مداخله‌ی بشردوستانه & \checkmark & &
  هسته‌ی اصلی پژوهش‌ها \\
۶. بازیابی/وحدت‌طلبی & بعضاً & بعضاً &
  بستگی به تعریف \lat{FIRC} دارد \\
۷. مداخله‌ی ضدتروریستی & \checkmark & &
  مثال‌های اصلی: عراق و افغانستان \\
۸. مداخله‌ی اقتصادی & & \checkmark &
  «نظامی» تلقی نمی‌شود \\
\bottomrule
\end{tabularx}
\caption*{\textsuperscript{*}مورد ژاپن و آلمان غربی گاه به‌عنوان
«نمونه‌ی موفق» مداخله ذکر می‌شوند — اما شرایط آنها
(شکست کامل نظامی، نهادهای موجود، حمایت مردمی، رقابت جنگ سرد)
به‌شدت استثنایی بود.}
\end{table}

\begin{tcolorbox}[enemybox, title={خطر خلط طبقه‌ها: دو مغالطه‌ی رایج}]

\textbf{مغالطه‌ی ۱: «ژاپن شد، پس عراق هم می‌شود»}

مدافعان حمله به عراق (۲۰۰۳) مکرراً به بازسازی موفق ژاپن اشاره کردند.
اما:
\begin{itemize}
\item ژاپن \textbf{طبقه‌ی ۲} (بازآرایی پساجنگی پس از شکست کامل) بود؛
  عراق \textbf{طبقه‌ی ۵/۷} (تغییر رژیم تحمیلی)
\item ژاپن \textbf{نهادهای بوروکراتیک} قوی داشت؛ عراق پس از
  «حل‌وفصل بعث» (\lat{De-Ba'athification}) نهادهایش نابود شد
\item ژاپن جامعه‌ی نسبتاً \textbf{همگن} بود؛ عراق جامعه‌ی
  \textbf{چندقومی-مذهبی} با تنش‌های عمیق
\item ژاپن در \textbf{شکست کامل} پذیرفته بود که باید تغییر کند؛
  عراقی‌ها چنین احساسی نداشتند
\item ژاپن از رقابت جنگ سرد سود برد (آمریکا انگیزه داشت آن را
  «ویترین» کند)؛ عراق چنین موقعیتی نداشت
\end{itemize}

\tcblower

\textbf{مغالطه‌ی ۲: «همه‌ی دخالت‌ها شکست‌اند، پس مَندیت هم شکست بود»}

منتقدان مداخله گاه مَندیت‌ها را در لیست «شکست مداخلات» می‌آورند. اما:
\begin{itemize}
\item مَندیت‌ها «مداخله‌ی رهایی‌بخش» نبودند — \keyword{ادعای رهایی}
  نداشتند
\item مشکل مَندیت‌ها نه «شکست مداخله» بلکه \keyword{ماهیت استعماری}
  خود مَندیت بود
\item خلط اینها، هم تحلیل استعمار را ضعیف می‌کند هم تحلیل مداخله را
\end{itemize}

\textbf{نتیجه:} دقت مفهومی نه فقط تمرین آکادمیک بلکه
\concept{ضرورت عملی} است. بدون تمایز دقیق، نه نقد مداخله ممکن
است نه دفاع از آن.
\end{tcolorbox}


%% ================================================================
\section{طیف مشروعیت: از سلطه تا همبستگی}
\label{sec:tax-legitimacy-spectrum}
%% ================================================================

\needspace{8\baselineskip}
اگر بخواهیم این هشت طبقه را بر یک \concept{طیف مشروعیت} قرار دهیم —
از کمترین مشروعیت (سلطه‌ی محض) تا بیشترین مشروعیت (همبستگی واقعی) — نتیجه
چنین خواهد بود:

\begin{tcolorbox}[timelinebox, title={طیف مشروعیت دخالت خارجی}]
\begin{table}[H]
\centering\small
\begin{tabularx}{\textwidth}{c|r|c|X}
\toprule
& \textbf{نوع} & \textbf{مشروعیت} & \textbf{توضیح} \\
\midrule
\enemy{۱} & کشورگشایی/استعمار & بسیار پایین &
  سلطه‌ی محض؛ امروز تقریباً هیچ توجیهی ندارد \\
\enemy{۲} & مداخله‌ی ایدئولوژیک مخفی & بسیار پایین &
  نقض حاکمیت بدون حتی اعلام عمومی \\
\enemy{۳} & بازیابی/وحدت‌طلبی تهاجمی & پایین &
  ناسیونالیسم پوشش تجاوز \\
\keyword{۴} & بازآرایی پساجنگی & مورد مناقشه &
  گاه «ضروری» (ژاپن)، گاه «استعمار نو» (مَندیت) \\
\keyword{۵} & مداخله‌ی ضدتروریستی & مورد مناقشه &
  تعریف «ترور» سیاسی است \\
\keyword{۶} & استعمارزدایی & مورد مناقشه &
  بستگی به نحوه‌ی خروج دارد \\
\keyword{۷} & مداخله‌ی بشردوستانه & متوسط رو به بالا &
  بالاترین ادعای اخلاقی — اما پرمخاطره‌ترین \\
\concept{۸} & مداخله‌ی اقتصادی ساختاری & مورد مناقشه &
  «کمک» یا «سلطه»؟ بستگی به شرایط دارد \\
\midrule
\multicolumn{4}{c}{\textit{فراتر از «مداخله»:}} \\
\midrule
\concept{۹} & همبستگی انتقادی & بالا &
  حمایت بدون سلطه (فصل~\ref{ch:responsibility-beyond-borders}) \\
\bottomrule
\end{tabularx}
\end{table}
\end{tcolorbox}


%% ================================================================
\section{جمع‌بندی: چرا تمایز مهم است}
\label{sec:tax-conclusion}
%% ================================================================

\begin{tcolorbox}[mybox, title={خلاصه‌ی این پیوست}]
\begin{enumerate}
\item «مداخله‌ی خارجی» اصطلاحی چتری است که پدیده‌های
  \textbf{بسیار متفاوتی} را پوشش می‌دهد — از استعمار تا مداخله‌ی
  بشردوستانه

\item حداقل \textbf{هشت طبقه‌ی متمایز} قابل شناسایی است:
  کشورگشایی، بازآرایی پساجنگی، استعمارزدایی، مداخله‌ی ایدئولوژیک،
  مداخله‌ی بشردوستانه، بازیابی/وحدت‌طلبی، مداخله‌ی ضدتروریستی،
  و مداخله‌ی اقتصادی ساختاری

\item \textbf{نظام مَندیت عثمانی} «مداخله‌ی خارجی» به‌معنای مدرن نبود
  — بلکه \concept{بازآرایی امپریالیستی پساجنگی} بود

\item \textbf{تقسیم هند} نتیجه‌ی «مداخله‌ی رهایی‌بخش» نبود — بلکه
  نتیجه‌ی \concept{خروج بد مدیریت‌شده‌ی استعمارگر} بود

\item \textbf{فلسطین} پیچیده‌ترین مورد است: تلاقی حداقل
  \concept{چهار طبقه} هم‌زمان

\item آمار «۸۴٪ شکست» مشخصاً به طبقه‌های ۴، ۵، و ۷ مربوط است —
  \textbf{خلط طبقه‌ها} هم به نقد مداخله آسیب می‌زند هم به دفاع از آن

\item این طبقه‌ها مستقل نیستند — \concept{زنجیره‌ای علّی} آنها را
  به‌هم پیوند می‌دهد: استعمار دیروز، زمینه‌ساز مداخله‌ی امروز است

\item فراتر از همه‌ی این طبقه‌ها، امکان \concept{همبستگی انتقادی}
  وجود دارد — شکلی از مسئولیت بین‌المللی که نه سلطه است نه بی‌تفاوتی
  (فصل~\ref{ch:responsibility-beyond-borders})
\end{enumerate}

\textbf{پیام نهایی:} دقت مفهومی فقط تمرین آکادمیک نیست — شرط لازم
\concept{تصمیم‌گیری اخلاقی} است. تا وقتی نمی‌دانیم دقیقاً با چه
پدیده‌ای روبه‌رو هستیم، نه می‌توانیم درست نقد کنیم، نه درست عمل.
\end{tcolorbox}

%% ────────────────────────────────────
\section*{منابع اصلی این پیوست}
\addcontentsline{toc}{section}{منابع اصلی این پیوست}
%% ────────────────────────────────────

\begin{itemize}[label={\color{PrimaryMid}\textbullet}]
\item Rosenau, James N. ``Intervention as a Scientific Concept.''
  \textit{Journal of Conflict Resolution} 13, no.\,2 (1969): 149--171.

\item Finnemore, Martha. \textit{The Purpose of Intervention: Changing
  Beliefs about the Use of Force}. Cornell UP, 2003.

\item Downes, Alexander B. \textit{Catastrophic Success: Why
  Foreign-Imposed Regime Change Goes Wrong}. Cornell UP, 2021.

\item O'Rourke, Lindsey A. \textit{Covert Regime Change: America's
  Secret Cold War}. Cornell UP, 2018.

\item Fromkin, David. \textit{A Peace to End All Peace: The Fall of the
  Ottoman Empire and the Creation of the Modern Middle East}. Holt, 1989.

\item Khalidi, Rashid. \textit{The Iron Cage: The Story of the
  Palestinian Struggle for Statehood}. Beacon Press, 2006.

\item Rogan, Eugene. \textit{The Fall of the Ottomans: The Great War in
  the Middle East}. Basic Books, 2015.

\item Pedersen, Susan. \textit{The Guardians: The League of Nations and
  the Crisis of Empire}. Oxford UP, 2015.

\item Khan, Yasmin. \textit{The Great Partition: The Making of India and
  Pakistan}. Yale UP, 2007.

\item Klein, Naomi. \textit{The Shock Doctrine: The Rise of Disaster
  Capitalism}. Picador, 2007.

\item Singer, Peter. ``Famine, Affluence, and Morality.''
  \textit{Philosophy \& Public Affairs} 1, no.\,3 (1972): 229--243.

\item Mamdani, Mahmood. \textit{Neither Settler nor Native: The Making
  and Unmaking of Permanent Minorities}. Harvard UP, 2020.

\item Fanon, Frantz. \textit{The Wretched of the Earth}. Trans.\
  Richard Philcox. Grove Press, 2004 [1961].
\end{itemize}

\cleardoublepage